\chapter{Introduction}\label{ch:intro}

Road vehicles are becoming networked computers. Through dedicated short-range radio, a vehicle can
tell the vehicles around it where it is, how fast it is going and that it is braking, long before a
driver could see the danger. Intelligent Transportation Systems (ITS) build collision warnings,
cooperative intersection management and emergency-vehicle preemption on these messages. Their value
depends on one question that the radio alone cannot answer: can the message be trusted? This report
describes a system that answers it without giving up the privacy of the people in the vehicles, and
measures what that costs on a realistic network.

\section{Vehicular Ad Hoc Networks}

A vehicular ad hoc network (VANET) consists of three kinds of participants:
\begin{itemize}
  \item \textbf{On-board units (OBUs):} the radio computers in vehicles. They broadcast a
  \emph{Basic Safety Message} (BSM)~\cite{saej2735}, also called a cooperative awareness message
  in Europe~\cite{etsi302637}, about ten times per second, with position, speed and heading.
  \item \textbf{Roadside units (RSUs):} fixed radios at junctions with a wired backbone connection.
  They authenticate vehicles, relay information and act as the network's local infrastructure.
  \item \textbf{A trusted authority (TA):} for example a transport department, which registers
  vehicles. It is not on the communication path during normal operation.
\end{itemize}
Vehicles talk to each other (V2V) and to RSUs (V2I) over IEEE 802.11p~\cite{jiang2008wave,kenney2011dsrc},
a Wi-Fi variant for the 5.9\,GHz band that works without association, so a vehicle can broadcast to
whoever is in range. Links last seconds, topology changes constantly, and every vehicle in range shares one channel.

\section{Motivation}

Three requirements pull in different directions:
\begin{enumerate}
  \item \textbf{Trust.} A forged BSM reporting a phantom vehicle can trigger an emergency brake; a
  forged emergency-vehicle message can clear an intersection for an attacker. Receivers must verify
  that every message comes from a registered vehicle, and misbehaving vehicles must be removed.
  \item \textbf{Privacy.} A vehicle that broadcasts its position ten times a second under a fixed
  identifier can be tracked through a city by anyone with a receiver~\cite{gruteser2005tracking}.
  Authentication must therefore not reveal \emph{which} vehicle is speaking, yet an authorized
  investigation must still be able to identify a vehicle that attacked the network (conditional
  privacy).
  \item \textbf{Efficiency.} Security bytes are paid ten times a second by every vehicle on a shared
  channel. Above roughly 60\,\% channel occupancy, broadcasts start to collide and safety messages
  are lost. Verification must also finish within milliseconds on modest automotive hardware.
\end{enumerate}
Certificate-based schemes~\cite{raya2007securing,brecht2018scms} solve the first by keeping a
central authority on the trust path and the second by changing pseudonym certificates, at the cost of
large certificate pools and revocation lists. Zero-knowledge schemes~\cite{rasheed2020adaptive,xi2022zama}
hide the credential but lack a shared, auditable record of authorization and revocation across RSUs.
Blockchain-based schemes~\cite{alsaeed2025baaiov,li2024bpa,wang2021btsca} provide that shared state
but add consensus latency. The Phase~I review of this project identified the open question:

\begin{quote}
\emph{Can privacy-preserving ZKP authentication be combined with permissioned blockchain
authorization and lifecycle management without unacceptable latency, computation or communication
overhead?}
\end{quote}

\section{Objectives}

\begin{enumerate}
  \item Design a scheme in which a vehicle proves registration with a zero-knowledge proof, an RSU
  decides access with a ledger policy, and registration roots, authorizations, misbehaviour reports,
  revocations and identity-reveal approvals live on a permissioned blockchain.
  \item Support the whole lifecycle: registration, authentication, handover, renewal, signed safety
  messages, detection and reporting, revocation and threshold identity reveal.
  \item Implement it completely: protocol, chaincode and gateway on Hyperledger Fabric, detector, and
  a network simulation on ns-3 with SUMO mobility.
  \item Measure it on realistic traffic against established baselines, report means with confidence
  intervals and significance tests, and state the limitations.
  \item Reduce the measured latency and overhead without weakening the security model, and handle
  emergency vehicles as a privileged class.
\end{enumerate}

\section{Contributions}

\begin{itemize}
  \item A complete authentication and authorization lifecycle that keeps the blockchain off the
  critical path of safety messages: proofs are verified at the RSU, and the ledger holds only the
  state that has to be shared and audited (Chapter~\ref{ch:design}).
  \item Mechanisms that make the scheme cheap on the radio channel and robust against misbehaving
  vehicles and RSUs: TESLA, implicit certificates, RSU clusters, dead reckoning, revocation filters,
  fast quarantine, trust-weighted reports and batch verification (Chapter~\ref{ch:opt}).
  \item An emergency-vehicle extension: class-bound credentials, a policy-restricted service,
  signed siren alerts that cannot be forged under a normal ticket, and a priority lane at RSUs.
  \item A reproducible evaluation framework: ns-3 802.11p with real Luxembourg traffic, measured
  cryptographic costs on a per-node processor model, the VeReMi Extension detection benchmark, a
  Fabric consensus benchmark, ten parameter sweeps, and a live browser visualisation for
  demonstrations (Chapters~\ref{ch:impl} and~\ref{ch:results}).
\end{itemize}

\section{Team and Division of Work}

The implementation was divided into four parts with written interfaces between them
(\texttt{docs/interfaces.md} in the repository). \textbf{Ben Savio} built the SUMO mobility and LuST
scenario, the ns-3 802.11p network, the ledger bridge and the integration of all parts;
\textbf{Harsh Kumar Gour} the Hyperledger Fabric network, the chaincode and the REST gateway;
\textbf{Byreddy Varshith Reddy} the authentication protocol (zero-knowledge login, tickets, handover
and revocation); and \textbf{Andela Harsha Vardhan} the misbehaviour detection and the evaluation
tooling.

\section{Report Structure}

Chapter~\ref{ch:background} covers the building blocks and the threat model, Chapter~\ref{ch:lit}
related work, Chapters~\ref{ch:design}--\ref{ch:impl} the scheme, its optimisations and the
implementation, Chapter~\ref{ch:results} the evaluation and Chapter~\ref{ch:conclusion} the
conclusions. The appendix explains how to reproduce every result.
